\documentclass[12pt, a4paper]{article}
\input{/home/mykhailo/.config/nvim/templates/preamble.tex}

\newcommand{\department}{ЗІ}
\newcommand{\reportnumber}{2}
\newcommand{\reporttitle}{Розробка консольних утиліт для задач кібербезпеки}
\newcommand{\profname}{Кебус В. М.}

\begin{document}

\input{/home/mykhailo/Documents/lpnu/report_title_page}
Мета:
\begin{enumerate}
	\item Закріпити та поглибити знання з об'єктно-орієнтованого програмування (ООП) в Python: створення класів, інкапсуляція, наслідування, композиція, спеціальні методи класів та декоратори \verb|@property|, \verb|@dataclass|.
	\item Опанувати роботу зі стандартною бібліотекою Python для обробки даних. Маніпуляція шляхами та файлами \verb|pathlib|, робота з форматами JSON та CSV, обробка дати і часу \verb|datetime|, використання спеціалізованих структур даних \verb|collections| та регулярних виразів \verb|re|.
	\item Навчитися будувати надійні та стійкі консольні програми (CLI-утиліти): обробка винятків, валідація вхідних даних, структуроване логування подій (\verb|logging|) та розбір аргументів командного рядка (\verb|argparse|).
\end{enumerate}
\end{document}
